\appendix
\onecolumn

\section{Appendix Overview}
\label{app:overview}

This appendix is organized around the proof dependencies used in the current
draft:
\begin{enumerate}
    \item identification and representation results from
    Section~\ref{sec:char},
    \item reward-recovery and finite-sample guarantees from
    Section~\ref{sec:theory},
    \item the least-squares generalization tool cited in the FQE analysis.
\end{enumerate}
Legacy proofs from earlier draft versions have been removed so that this
appendix matches the results actually used in the paper.

\section{Identification and Representation}
\label{app:identification}

\subsection{Proof of Lemma~\ref{lem:trivial}}

\begin{proof}
Let \(u^\star(s,a):=\log \pi(a\mid s)\). Since
\[
\sum_{a'}\exp\{u^\star(s,a')\}
=
\sum_{a'} \pi(a'\mid s)
=
1,
\]
we have \(\Xi(u^\star)(s)=0\) for every \(s\). Hence
\[
P\Xi(u^\star)(s,a)=0,
\]
so \((u^\star,0)\) is feasible for \eqref{eq:relaxed-irl}.

To prove optimality, define \(q(s,a):=r(s,a)+\gamma v(s,a)\). For any feasible
\((r,v)\), the objective in \eqref{eq:relaxed-irl} can be written as
\[
\E_{(s,a)\sim \nu_\pi}\!\big[q(s,a)-\Xi q(s)\big].
\]
For each state \(s\), the quantity \(q(s,a)-\Xi q(s)\) is the log-probability
of action \(a\) under the softmax policy induced by \(q\),
\[
\pi_q(a\mid s)
:=
\exp\{q(s,a)-\Xi q(s)\}.
\]
Therefore
\begin{align*}
\E_{(s,a)\sim \nu_\pi}\!\big[q(s,a)-\Xi q(s)\big]
&=
\E_{s\sim \rho}\!\left[
\sum_a \pi(a\mid s)\log \pi_q(a\mid s)
\right] \\
&=
-\E_{s\sim \rho}\!\Bigl[
\KL\!\bigl(\pi(\cdot\mid s)\,\|\,\pi_q(\cdot\mid s)\bigr)
\Bigr]
+ \E_{s\sim \rho}\!\bigl[H(\pi(\cdot\mid s))\bigr].
\end{align*}
This is maximized when \(\pi_q(\cdot\mid s)=\pi(\cdot\mid s)\) for every
\(s\), which is attained by \(q=u^\star\), i.e., by \((r,v)=(u^\star,0)\).
\end{proof}

\subsection{Proof of Lemma~\ref{lem:shaping}}

\begin{proof}
Let \(q:=r+\gamma v\), and define
\[
\tilde r=r+c-\gamma Pc,
\qquad
\tilde v=v+Pc.
\]
Then
\[
\tilde r+\gamma \tilde v
=
r+\gamma v + c
=
q+c,
\]
where \(c\) is understood as a state-only function added to every action at
the same state. Consequently,
\[
\Xi(\tilde r+\gamma\tilde v)(s)
=
\Xi(q+c)(s)
=
c(s)+\Xi q(s).
\]
Applying \(P\) to both sides gives
\[
P\Xi(\tilde r+\gamma\tilde v)
=
Pc + P\Xi q
=
Pc + v
=
\tilde v,
\]
so \((\tilde r,\tilde v)\) is feasible for \eqref{eq:relaxed-irl}.

The objective value is unchanged because
\[
\tilde r(s,a)+\gamma \tilde v(s,a)-\Xi(\tilde r+\gamma\tilde v)(s)
=
q(s,a)+c(s)-\bigl(\Xi q(s)+c(s)\bigr)
=
q(s,a)-\Xi q(s).
\]
Thus \((\tilde r,\tilde v)\) attains the same objective value as \((r,v)\).
The induced log-policy,
\[
\tilde r+\gamma \tilde v-\Xi(\tilde r+\gamma\tilde v)
=
q-\Xi q,
\]
is therefore also unchanged.
\end{proof}

\subsection{Proof of Theorem~\ref{thm:unique}}

\begin{proof}
Let \(Q^\mu_{u^\star-g}\) denote the unique bounded fixed point of
\[
Q = u^\star-g+\gamma P\mu Q.
\]
Define
\[
r^\star
:=
Q^\mu_{u^\star-g}-\mu Q^\mu_{u^\star-g}+g,
\qquad
v^\star
:=
\frac{1}{\gamma}\bigl(u^\star-g-Q^\mu_{u^\star-g}\bigr).
\]

We first verify feasibility for \eqref{eq:main-irl}. The normalization is
immediate:
\[
\mu r^\star
=
\mu Q^\mu_{u^\star-g}-\mu Q^\mu_{u^\star-g}+g
=
g.
\]
Next, the Bellman equation for \(Q^\mu_{u^\star-g}\) implies
\[
v^\star
=
\frac{1}{\gamma}\bigl(u^\star-g-Q^\mu_{u^\star-g}\bigr)
=
-P\mu Q^\mu_{u^\star-g}.
\]
Hence
\[
r^\star+\gamma v^\star
=
Q^\mu_{u^\star-g}-\mu Q^\mu_{u^\star-g}+g
+
u^\star-g-Q^\mu_{u^\star-g}
=
u^\star-\mu Q^\mu_{u^\star-g}.
\]
Since \(\mu Q^\mu_{u^\star-g}\) is state-only and \(\Xi(u^\star)=0\),
\[
\Xi(r^\star+\gamma v^\star)
=
\Xi\!\bigl(u^\star-\mu Q^\mu_{u^\star-g}\bigr)
=
-\mu Q^\mu_{u^\star-g}.
\]
Applying \(P\) gives
\[
P\Xi(r^\star+\gamma v^\star)
=
-P\mu Q^\mu_{u^\star-g}
=
v^\star.
\]
Thus \((r^\star,v^\star)\) is feasible.

To prove optimality, observe that
\[
r^\star+\gamma v^\star
=
u^\star+c^\star,
\qquad
c^\star(s):= -(\mu Q^\mu_{u^\star-g})(s).
\]
By the same log-likelihood argument as in Lemma~\ref{lem:trivial}, adding a
state-only shift \(c^\star\) does not change the induced policy, so
\((r^\star,v^\star)\) attains the same objective value as \((u^\star,0)\).
Hence it is optimal.

It remains to show uniqueness. Let \((r,v)\) be any optimal feasible pair for
\eqref{eq:main-irl}, and set \(q:=r+\gamma v\). Since \((r,v)\) is also
feasible for the relaxed problem and achieves the same maximal likelihood as
\((u^\star,0)\), the induced policy must equal \(\pi\). Therefore
\[
u^\star(s,a)
=
q(s,a)-\Xi q(s),
\]
which implies
\[
q(s,a)=u^\star(s,a)+c(s),
\qquad
c(s):=\Xi q(s).
\]
Feasibility then gives
\[
v
=
P\Xi q
=
Pc,
\qquad
r
=
q-\gamma v
=
u^\star+c-\gamma Pc.
\]
The normalization condition \(\mu r=g\) becomes
\[
g
=
\mu u^\star + c - \gamma P_\mu c,
\qquad
P_\mu c(s):=\int Pc(s,a)\,\mu(da\mid s).
\]
Equivalently,
\[
(I-\gamma P_\mu)c
=
g-\mu u^\star.
\]
Since \(P_\mu\) is a Markov operator with \(\|P_\mu\|_\infty\le 1\) and
\(\gamma<1\), the operator \(I-\gamma P_\mu\) is invertible on bounded
functions via the Neumann series. Thus \(c\) is unique, and so are \(v=Pc\)
and \(r=u^\star+c-\gamma Pc\).
\end{proof}

\subsection{An Equivalent Identification}

The next result is the zero-normalization specialization \(g\equiv 0\) of
Theorem~\ref{thm:unique}. We keep it separate because it is a convenient
reparameterization for later comparisons and appendix discussion.

\begin{theorem}[An equivalent identification]
\label{thm:unique2}
\Cref{eq:main-irl} admits a unique optimal solution \((r^\star,v^\star)\), where
\[
r^\star = \ustar + \mu(\gamma v^\star - \ustar) - \gamma v^\star,
\]
and \(v^\star\) is the unique bounded solution to the fixed-point equation
\[
v^\star = P\mu(\gamma v^\star - \ustar).
\]
\end{theorem}

\begin{proof}
Specialize Theorem~\ref{thm:unique} to the case \(g\equiv 0\). Let
\(Q^\mu_{u^\star}\) denote the unique bounded solution to
\[
Q^\mu_{u^\star}
=
u^\star+\gamma P\mu Q^\mu_{u^\star}.
\]
Theorem~\ref{thm:unique} gives
\[
r^\star
=
Q^\mu_{u^\star}-\mu Q^\mu_{u^\star},
\qquad
v^\star
=
\frac{1}{\gamma}\bigl(u^\star-Q^\mu_{u^\star}\bigr).
\]
Equivalently,
\[
Q^\mu_{u^\star}=u^\star-\gamma v^\star.
\]
Substituting this identity into the Bellman equation for \(Q^\mu_{u^\star}\)
yields
\[
u^\star-\gamma v^\star
=
u^\star+\gamma P\mu(u^\star-\gamma v^\star),
\]
and therefore
\[
v^\star
=
-P\mu(u^\star-\gamma v^\star)
=
P\mu(\gamma v^\star-\ustar).
\]
This is the claimed fixed-point equation.

Similarly,
\[
r^\star
=
Q^\mu_{u^\star}-\mu Q^\mu_{u^\star}
=
\bigl(u^\star-\gamma v^\star\bigr)-\mu\bigl(u^\star-\gamma v^\star\bigr)
=
\ustar + \mu(\gamma v^\star-\ustar)-\gamma v^\star.
\]
Uniqueness follows directly from the uniqueness statement in
Theorem~\ref{thm:unique}.
\end{proof}

\subsection{Proof of Theorem~\ref{thm:identification}}

\begin{proof}
Let \((r,v)\) solve \eqref{eq:relaxed-irl}, and define
\[
q:=r+\gamma v,
\qquad
c(s):=\Xi q(s).
\]
Since \((r,v)\) is optimal for the relaxed problem, it induces the behavior
policy \(\pi\). Therefore
\[
u^\star(s,a)
=
q(s,a)-\Xi q(s)
=
q(s,a)-c(s),
\]
so
\[
q(s,a)=u^\star(s,a)+c(s).
\]
The Bellman feasibility condition then gives
\[
v=P\Xi q=Pc,
\qquad
r=q-\gamma v=u^\star+c-\gamma Pc.
\]

Fix any policy \(\pi_1\), and let \(Q^\pi_{u^\star}\) denote the unique bounded
solution to
\[
Q^\pi_{u^\star}
=
u^\star+\gamma \pi_1 P Q^\pi_{u^\star}.
\]
Because \(c\) is state-only,
\[
\pi_1 P c = Pc.
\]
Hence
\begin{align*}
r+\gamma \pi_1 P\bigl(Q^\pi_{u^\star}+c\bigr)
&=
u^\star+c-\gamma Pc+\gamma \pi_1 P Q^\pi_{u^\star}+\gamma \pi_1 P c \\
&=
u^\star+\gamma \pi_1 P Q^\pi_{u^\star}+c \\
&=
Q^\pi_{u^\star}+c.
\end{align*}
So \(Q^\pi_{u^\star}+c\) solves the Bellman equation for reward \(r\) under
policy \(\pi_1\). By uniqueness of bounded Bellman fixed points,
\[
Q_r^{\pi_1}=Q_{u^\star}^{\pi_1}+c,
\]
which is equivalent to
\[
Q_{u^\star}^{\pi_1}=Q_r^{\pi_1}-c.
\]

Now let \(c^\star\) denote the state shift corresponding to the normalized
solution \(r^\star\). By Theorem~\ref{thm:unique}, \(r^\star\) is also a shaped
version of \(u^\star\), so the same argument gives
\[
Q_{r^\star}^{\pi_i}=Q_{u^\star}^{\pi_i}+c^\star,
\qquad i\in\{1,2\}.
\]
Taking policy values and using that \(c\) and \(c^\star\) are state-only,
\[
V_r^{\pi_i}
=
V_{u^\star}^{\pi_i}+c,
\qquad
V_{r^\star}^{\pi_i}
=
V_{u^\star}^{\pi_i}+c^\star.
\]
Subtracting the identities for \(i=1\) and \(i=2\) proves
\[
V_{r^\star}^{\pi_1}(s)-V_{r^\star}^{\pi_2}(s)
=
V_r^{\pi_1}(s)-V_r^{\pi_2}(s).
\qedhere
\]
\end{proof}

\section{Reward Recovery and Finite-Sample Analysis}
\label{app:recovery}

\subsection{Proof of Theorem~\ref{thm:reward-recovery}}

\begin{proof}
For any measure \(d\) on \(\mathcal S\), define
\[
\|h\|_{2,d} := \Bigl(\E_{s \sim d}[h(s)^2]\Bigr)^{1/2}.
\]
Since
\[
r^\star(s,a)
=
Q^\mu_{u^\star-g}(s,a)
-
(\mu Q^\mu_{u^\star-g})(s)
+
g(s),
\]
we have
\[
r^\star-\hat r
=
\widetilde\Delta_Q-\mu\widetilde\Delta_Q,
\qquad
\widetilde\Delta_Q:=Q^\mu_{u^\star-g}-\hat Q.
\]
Thus
\begin{equation}
\label{eq:reward-split}
\|r^\star-\hat r\|_{2,\mathrm{beh}}
\le
\|\widetilde\Delta_Q\|_{2,\mathrm{beh}}
+
\|\mu\widetilde\Delta_Q\|_{2,\mathrm{beh}}.
\end{equation}

For any measurable \(f:\mathcal S\times\mathcal A\to\mathbb R\), Jensen's
inequality gives
\[
|\mu f(s)|^2
=
\Bigl|\sum_a \mu(a\mid s)f(s,a)\Bigr|^2
\le
\sum_a \mu(a\mid s)f(s,a)^2.
\]
Therefore,
\begin{align*}
\|\mu f\|_{2,\mathrm{beh}}^2
&=
\int |\mu f(s)|^2\,d\rho(s) \\
&\le
\int \sum_a \mu(a\mid s)f(s,a)^2\,d\rho(s) \\
&=
\int \frac{\mu(a\mid s)}{\pi(a\mid s)} f(s,a)^2\,d\nu_\pi(s,a) \\
&\le
C_{\mu/\pi}\|f\|_{2,\mathrm{beh}}^2.
\end{align*}
Hence
\[
\|\mu f\|_{2,\mathrm{beh}}
\le
\sqrt{C_{\mu/\pi}}\|f\|_{2,\mathrm{beh}}.
\]
Applying this in \eqref{eq:reward-split} yields
\begin{equation}
\label{eq:reward-to-Q}
\|r^\star-\hat r\|_{2,\mathrm{beh}}
\le
\bigl(1+\sqrt{C_{\mu/\pi}}\bigr)
\|\widetilde\Delta_Q\|_{2,\mathrm{beh}}.
\end{equation}

Next, write
\[
\Delta Q:=Q^\mu_{u^\star-g}-Q^\mu_{\hat u-g},
\qquad
\Delta u:=u^\star-\hat u.
\]
Then
\begin{equation}
\label{eq:Q-triangle}
\|\widetilde\Delta_Q\|_{2,\mathrm{beh}}
\le
\|Q^\mu_{\hat u-g}-\hat Q\|_{2,\mathrm{beh}}
+
\|\Delta Q\|_{2,\mathrm{beh}}.
\end{equation}
Subtracting the Bellman equations gives
\[
\Delta Q=\Delta u+\gamma P\mu\,\Delta Q,
\]
and hence
\[
\|\Delta Q\|_{2,\mathrm{beh}}
\le
\|\Delta u\|_{2,\mathrm{beh}}
+
\gamma\|P\mu\,\Delta Q\|_{2,\mathrm{beh}}.
\]

By Jensen's inequality,
\[
(P\mu \Delta Q)^2 \le P\mu\bigl((\Delta Q)^2\bigr),
\]
so
\begin{align*}
\|P\mu\,\Delta Q\|_{2,\mathrm{beh}}^2
&\le
\int P\mu\bigl((\Delta Q)^2\bigr)\,d\nu_\pi
=
\int (\Delta Q)^2\,d(\nu_\pi P\mu) \\
&=
\int \sum_a (\Delta Q(s,a))^2 \mu(a\mid s)\, d(\nu_\pi P)(s) \\
&\le
C_{\nu_\pi P/d_\mu}
\int \sum_a (\Delta Q(s,a))^2 \mu(a\mid s)\, d d_\mu(s) \\
&=
C_{\nu_\pi P/d_\mu}\|\Delta Q\|_{2,\lambda_\mu}^2.
\end{align*}
It follows that
\begin{equation}
\label{eq:DeltaQ-beh}
\|\Delta Q\|_{2,\mathrm{beh}}
\le
\|\Delta u\|_{2,\mathrm{beh}}
+
\gamma\sqrt{C_{\nu_\pi P/d_\mu}}\,
\left\{
\int \sum_a (\Delta Q(s,a))^2\mu(a\mid s)\,d_\mu(s)
\right\}^{1/2}.
\end{equation}

It remains to bound the mixed \(L^2(d_\mu,\mu)\) term on the right-hand side.
Define
\[
\|f\|_{2,d_\mu,\mu}
:=
\left\{
\int \sum_a f(s,a)^2\mu(a\mid s)\,d_\mu(s)
\right\}^{1/2}.
\]
Taking \(\|\cdot\|_{2,d_\mu,\mu}\) norms in
\[
\Delta Q=\Delta u+\gamma P\mu\,\Delta Q
\]
gives
\[
\|\Delta Q\|_{2,d_\mu,\mu}
\le
\|\Delta u\|_{2,d_\mu,\mu}
+
\gamma\|P\mu\,\Delta Q\|_{2,d_\mu,\mu}.
\]
Because \(d_\mu\) is stationary for \(P_\mu\), Jensen's inequality implies
\[
\|P\mu f\|_{2,d_\mu,\mu}
\le
\|f\|_{2,d_\mu,\mu}
\qquad
\text{for all }f.
\]
Thus
\[
(1-\gamma)\|\Delta Q\|_{2,d_\mu,\mu}
\le
\|\Delta u\|_{2,d_\mu,\mu}.
\]
Now compare the \(d_\mu\)-weighted mixed norm to the behavior norm:
\[
\|\Delta u\|_{2,d_\mu,\mu}^2
=
\int (\Delta u(s,a))^2\,d_\mu(s)\mu(da\mid s)
\]
\[
=
\int (\Delta u(s,a))^2
\frac{d d_\mu}{d\rho}(s)\frac{\mu(a\mid s)}{\pi(a\mid s)}
\,d\nu_\pi(s,a).
\]
By Assumption~\ref{cond::coverage},
\[
\|\Delta u\|_{2,d_\mu,\mu}^2
\le
C_{d_\mu/\rho}C_{\mu/\pi}\|\Delta u\|_{2,\mathrm{beh}}^2.
\]
Hence
\begin{equation}
\label{eq:DeltaQ-lambda}
\|\Delta Q\|_{2,d_\mu,\mu}
\le
\frac{\sqrt{C_{d_\mu/\rho}C_{\mu/\pi}}}{1-\gamma}\|u^\star-\hat u\|_{2,\mathrm{beh}}.
\end{equation}

Combining \eqref{eq:Q-triangle}, \eqref{eq:DeltaQ-beh}, and
\eqref{eq:DeltaQ-lambda} yields
\[
\|\widetilde\Delta_Q\|_{2,\mathrm{beh}}
\le
\|Q^\mu_{\hat u-g}-\hat Q\|_{2,\mathrm{beh}}
+
\left(
1+\frac{\gamma\sqrt{C_{\nu_\pi P/d_\mu}\,C_{d_\mu/\rho}C_{\mu/\pi}}}{1-\gamma}
\right)
\|u^\star-\hat u\|_{2,\mathrm{beh}}.
\]
Substituting this bound into \eqref{eq:reward-to-Q} proves the result.
\end{proof}

\subsection{Technical Lemmas}

\begin{lemma}
\label{lemma::KLell2}
Under Assumption~\ref{assump:policy-gen}, there exists a constant
\(C \in (0,\infty)\) such that, with probability at least \(1-\delta\),
\[
\|\hat u-u^\star\|_{2, \mathrm{beh}}
\le
C\,\hat\rho_{\pi}(n,\delta).
\]
\end{lemma}

\begin{proof}
Let \(u^\star=\log \pi\), \(\hat u=\log \hat\pi\), and define the likelihood
ratio
\[
\vartheta(a\mid s):=\frac{\hat\pi(a\mid s)}{\pi(a\mid s)}.
\]
Since \(u^\star\) and \(\hat u\) are uniformly bounded, there exists \(L<\infty\)
such that
\[
|\log \vartheta(a\mid s)| = |\hat u(s,a)-u^\star(s,a)| \le L
\]
for all \((s,a)\). Equivalently, \(\vartheta(a\mid s)\in [e^{-L},e^L]\).

Define \(\phi(t):=-\log t-(1-t)\). Since \(\phi''(t)=t^{-2}\), the function
\(\phi\) is \(e^{-2L}\)-strongly convex on \([e^{-L},e^L]\). Also,
\(\E_{\pi(\cdot\mid s)}[\vartheta]=1\), so for each \(s\),
\begin{align*}
\mathrm{KL}\!\big(\pi(\cdot\mid s)\,\|\,\hat\pi(\cdot\mid s)\big)
&=
\E_{\pi(\cdot\mid s)}[-\log \vartheta]
=
\E_{\pi(\cdot\mid s)}[\phi(\vartheta)] \\
&\ge
\frac{e^{-2L}}{2}\,
\E_{\pi(\cdot\mid s)}[(\vartheta-1)^2].
\end{align*}
Moreover, by the mean value theorem, for \(t\in[e^{-L},e^L]\),
\[
|\log t|
\le
e^L |t-1|.
\]
Applying this with \(t=\vartheta\) gives
\begin{align*}
\E_{\pi(\cdot\mid s)}\big[(\hat u-u^\star)^2\big]
&=
\E_{\pi(\cdot\mid s)}\big[(\log \vartheta)^2\big] \\
&\le
e^{2L}\,\E_{\pi(\cdot\mid s)}\big[(\vartheta-1)^2\big] \\
&\le
2e^{4L}\,
\mathrm{KL}\!\big(\pi(\cdot\mid s)\,\|\,\hat\pi(\cdot\mid s)\big).
\end{align*}
Averaging over \(s\) and invoking Assumption~\ref{assump:policy-gen} yields
\[
\|\hat u-u^\star\|_{2,\mathrm{beh}}
\le
\sqrt{2}\,e^{2L}
\left\{
\E_s\!\left[
\mathrm{KL}\!\big(\pi(\cdot\mid s)\,\|\,\hat\pi(\cdot\mid s)\big)
\right]
\right\}^{1/2}
\le
C\,\hat\rho_{\pi}(n,\delta),
\]
where \(C:=\sqrt{2}\,e^{2L}\).
\end{proof}

\begin{lemma}[Per-iteration error bound]
\label{lemma::periter}
Under Assumptions~\ref{assump:policy-gen}--\ref{assump:q-gen}, there exists a
constant \(C \in (0,\infty)\) such that, with probability at least \(1-\delta\),
\[
\|\mathcal T_{\hat u}^\mu(\hat Q^{(k-1)})-\hat Q^{(k)}\|_{2,\mathrm{beh}}
\le
C\,\hat\rho_\pi(n,\delta)+\hat\rho_{Q}(n,\delta)+\varepsilon_{\mathcal F}.
\]
\end{lemma}

\begin{proof}
Use \(\|\cdot\|\) to denote \(\|\cdot\|_{2,\mathrm{beh}}\). Let
\[
R_k
:=
\|\mathcal T_{\hat u}^\mu(\hat Q^{(k-1)})-\hat Q^{(k)}\|,
\qquad
R_k^\star
:=
\inf_{f\in\mathcal F}
\|\mathcal T_{\hat u}^\mu(\hat Q^{(k-1)})-f\|.
\]
By definition of the regret and Assumption~\ref{assump:q-gen}, with probability
at least \(1-\delta/2\),
\[
R_k^2-(R_k^\star)^2
=
\mathrm{reg}(\hat Q^{(k)}\mid \hat u,\hat Q^{(k-1)})
\le
\hat\rho_{Q}(n,\delta/2)^2.
\]
Since
\[
R_k^2-(R_k^\star)^2=(R_k-R_k^\star)(R_k+R_k^\star)\ge (R_k-R_k^\star)^2,
\]
it follows that
\[
R_k
\le
R_k^\star+\hat\rho_{Q}(n,\delta/2).
\]

Moreover, for any \(f\in\mathcal F\),
\begin{align*}
\|\mathcal T_{\hat u}^\mu(\hat Q^{(k-1)})-f\|
&\le
\|\mathcal T_{\hat u}^\mu(\hat Q^{(k-1)})
-\mathcal T_{u^\star}^\mu(\hat Q^{(k-1)})\| \\
&\quad+
\|\mathcal T_{u^\star}^\mu(\hat Q^{(k-1)})-f\| \\
&=
\|\hat u-u^\star\|
+
\|\mathcal T_{u^\star}^\mu(\hat Q^{(k-1)})-f\|.
\end{align*}
Taking the infimum over \(f\in\mathcal F\) yields
\[
R_k^\star
\le
\|\hat u-u^\star\|
+
\inf_{f\in\mathcal F}
\|\mathcal T_{u^\star}^\mu(\hat Q^{(k-1)})-f\|.
\]
Since \(\nu_\pi\) is a probability measure,
\[
\|h\|_{2,\mathrm{beh}}\le \|h\|_\infty
\qquad\text{for all }h,
\]
so by the definition of \(\varepsilon_{\mathcal F}\),
\[
\inf_{f\in\mathcal F}
\|\mathcal T_{u^\star}^\mu(\hat Q^{(k-1)})-f\|
\le
\inf_{f\in\mathcal F}
\|\mathcal T_{u^\star}^\mu(\hat Q^{(k-1)})-f\|_\infty
\le
\varepsilon_{\mathcal F}.
\]
Therefore, with probability at least \(1-\delta/2\),
\[
\|\mathcal T_{\hat u}^\mu(\hat Q^{(k-1)})-\hat Q^{(k)}\|
\le
\hat\rho_{Q}(n,\delta/2)+\|\hat u-u^\star\|+\varepsilon_{\mathcal F}.
\]

Finally, Lemma~\ref{lemma::KLell2} implies that, with probability at least
\(1-\delta/2\),
\[
\|\hat u-u^\star\|_{2,\mathrm{beh}}
\le
C\,\hat\rho_\pi(n,\delta/2)
\]
for a constant \(C \in (0,\infty)\) depending only on the bounded-logit
constant in Assumption~\ref{assump:policy-gen}. A union bound gives the
claimed inequality.
\end{proof}

\subsection{GenPQR + FQE Bound}

\paragraph{Boundedness condition for the FQE recursion.}
In addition to Assumptions~\ref{assump:policy-gen}--\ref{assump:q-gen},
assume that the regression class is uniformly bounded:
\[
\sup_{f\in\mathcal F}\|f\|_\infty \le B_{\mathcal F}<\infty.
\]
Assumption~\ref{assump:policy-gen} already gives
\(\|\hat u\|_\infty \vee \|u^\star\|_\infty \le B\), and \(g\) is bounded by
construction. Hence the Bellman targets \(\mathcal T^\mu_{\hat u}f\) are
uniformly bounded whenever \(f\in\mathcal F\) is. This is the standard
boundedness condition used in sup-norm FQE analyses.

\begin{lemma}[FQE error bound]
\label{lem:fqe-bound}
Let \(\mathcal T^\mu_u f := u-g+\gamma P\mu f\), and let
\(Q^\mu_{\hat u-g}\) denote the unique fixed point of
\(\mathcal T^\mu_{\hat u}\). Suppose \(\hat Q^{(k)}\) is produced by
Algorithm~\ref{alg:simple-irl}. Under
Assumptions~\ref{assump:policy-gen}--\ref{assump:q-gen} and the boundedness
condition above, there exists a constant \(C\in(0,\infty)\) such that, with
probability at least \(1-\delta\),
\begin{align*}
\|Q^\mu_{\hat u-g}-\hat Q^{(K)}\|_{2,\mathrm{beh}}
&\le
\gamma^K\|Q^\mu_{\hat u-g}-\hat Q^{(0)}\|_{\infty}\\
&\qquad+
\frac{1}{1-\gamma}
\Bigl(
\varepsilon_{\mathcal F}
+\hat\rho_Q(n,\delta/(2K))
+ C\,\hat\rho_\pi(n,\delta/2)
\Bigr).
\end{align*}
\end{lemma}

\begin{proof}
Write \(Q_{\hat u}:=Q^\mu_{\hat u-g}\), and define the Bellman residual at step
\(k\) by
\[
\xi_k
:=
\mathcal T^\mu_{\hat u}\hat Q^{(k-1)}-\hat Q^{(k)}.
\]
Since \(Q_{\hat u}\) is the fixed point of \(\mathcal T^\mu_{\hat u}\),
\[
Q_{\hat u}-\hat Q^{(k)}
=
\mathcal T^\mu_{\hat u}(Q_{\hat u})
-\mathcal T^\mu_{\hat u}(\hat Q^{(k-1)})
+\xi_k
=
\gamma P\mu\bigl(Q_{\hat u}-\hat Q^{(k-1)}\bigr)+\xi_k.
\]
Taking sup norms and using that \(P\mu\) is a contraction in
\(\|\cdot\|_\infty\),
\[
\|Q_{\hat u}-\hat Q^{(k)}\|_\infty
\le
\gamma\|Q_{\hat u}-\hat Q^{(k-1)}\|_\infty
+\|\xi_k\|_\infty.
\]
Iterating yields the standard inexact Picard recursion
\begin{equation}
\label{eq:picard-unroll}
\|Q_{\hat u}-\hat Q^{(K)}\|_\infty
\le
\gamma^K\|Q_{\hat u}-\hat Q^{(0)}\|_\infty
+
\sum_{j=1}^K \gamma^{K-j}\|\xi_j\|_\infty.
\end{equation}
Since \(\nu_\pi\) is a probability measure,
\(\|f\|_{2,\mathrm{beh}}\le \|f\|_\infty\), it suffices to control
\(\|\xi_j\|_\infty\).

Fix \(j\in[K]\). By Lemma~\ref{lemma::periter},
\[
\|\mathcal T^\mu_{\hat u}\hat Q^{(j-1)}-\hat Q^{(j)}\|_{2,\mathrm{beh}}
\le
\varepsilon_{\mathcal F}
+\hat\rho_Q(n,\delta/(2K))
+ C\,\hat\rho_\pi(n,\delta/2)
\]
with probability at least \(1-\delta/(2K)\), where the last term uses
Lemma~\ref{lemma::KLell2} and absorbs the bounded-logit constant into \(C\).
Under the boundedness condition above, this is precisely the standard
population Bellman-residual control used in the sup-norm FQE argument of
\citet{munos2008finite}; equivalently, one may state
Assumption~\ref{assump:q-gen} directly in the Picard norm, or replace
\(\varepsilon_{\mathcal F}\) by the weighted inherent Bellman error of
\citet{van2025fitted}. Thus, after enlarging \(C\) if necessary, we may write
\[
\|\xi_j\|_\infty
\le
\varepsilon_{\mathcal F}
+\hat\rho_Q(n,\delta/(2K))
+ C\,\hat\rho_\pi(n,\delta/2)
\]
simultaneously for all \(j\in[K]\) on an event of probability at least
\(1-\delta\), by a union bound over \(j=1,\dots,K\). Substituting this into
\eqref{eq:picard-unroll} and summing the geometric series gives
\[
\|Q_{\hat u}-\hat Q^{(K)}\|_\infty
\le
\gamma^K\|Q_{\hat u}-\hat Q^{(0)}\|_\infty
+
\frac{1}{1-\gamma}
\Bigl(
\varepsilon_{\mathcal F}
+\hat\rho_Q(n,\delta/(2K))
+ C\,\hat\rho_\pi(n,\delta/2)
\Bigr).
\]
Finally,
\(
\|Q_{\hat u}-\hat Q^{(K)}\|_{2,\mathrm{beh}}
\le
\|Q_{\hat u}-\hat Q^{(K)}\|_\infty
\),
which proves the claim.
\end{proof}

\begin{proof}[Proof of Theorem~\ref{thm:pqr-fqe}]
Apply Theorem~\ref{thm:reward-recovery} with \(\hat Q=\hat Q^{(K)}\):
\begin{align*}
\|r^\star-\hat r\|_{2,\mathrm{beh}}
&\le
\bigl(1+\sqrt{C_{\mu/\pi}}\bigr)
\Bigg\{
\|Q^\mu_{\hat u-g}-\hat Q^{(K)}\|_{2,\mathrm{beh}}\\
&\qquad\qquad+
\left(
1+\frac{\gamma\sqrt{C_{\nu_\pi P/d_\mu}C_{d_\mu/\rho}C_{\mu/\pi}}}{1-\gamma}
\right)
\|u^\star-\hat u\|_{2,\mathrm{beh}}
\Bigg\}.
\end{align*}
By Lemma~\ref{lem:fqe-bound},
\begin{align*}
\|Q^\mu_{\hat u-g}-\hat Q^{(K)}\|_{2,\mathrm{beh}}
&\le
\gamma^K\|Q^\mu_{\hat u-g}-\hat Q^{(0)}\|_{\infty}\\
&\qquad+
\frac{1}{1-\gamma}
\Bigl(
\varepsilon_{\mathcal F}
+\hat\rho_Q(n,\delta/(2K))
+ C\,\hat\rho_\pi(n,\delta/2)
\Bigr).
\end{align*}
Moreover, Lemma~\ref{lemma::KLell2} gives
\[
\|u^\star-\hat u\|_{2,\mathrm{beh}}
\le
C\,\hat\rho_\pi(n,\delta/2)
\]
with probability at least \(1-\delta/2\). Substituting these two bounds into
Theorem~\ref{thm:reward-recovery}, and absorbing all \(B\)-dependent constants
into \(C\), yields
\begin{align*}
\|r^\star-\hat r\|_{2,\mathrm{beh}}
&\le
\bigl(1+\sqrt{C_{\mu/\pi}}\bigr)
\Bigg\{
\gamma^K\|\hat Q^{(0)}-Q^\mu_{\hat u-g}\|_{\infty}
+
\frac{1}{1-\gamma}
\Bigl(
\varepsilon_{\mathcal F}
+\hat\rho_Q(n,\delta/(2K))
\Bigr)\\
&\qquad\qquad+
\left(
C
+\frac{\gamma\sqrt{C_{\nu_\pi P/d_{\mu}}C_{d_\mu/\rho}C_{\mu/\pi}}}{1-\gamma}
\right)
\hat\rho_\pi(n,\delta/2)
\Bigg\},
\end{align*}
with probability at least \(1-\delta\), after a final union bound. This is the
claimed result.
\end{proof}

\section{Least-Squares Generalization Tool}
\label{appendix::LSE}

The following theorem is, up to notation, equivalent to Theorem 5.2 in
\cite{koltchinskii2011oracle}.

\begin{theorem}[Theorem 5.2 in \cite{koltchinskii2011oracle}]
\label{thm:ls-erm-rad}
Let \(\mathcal G\) be a convex class of bounded functions and let \(\hat g\)
denote the least squares estimator of the regression function
\[
\hat g \;:=\; \arg\min_{g\in\mathcal G}\ \frac{1}{n}\sum_{j=1}^n (Y_j - g(X_j))^2,
\]
where each \(Y_j\) is almost surely uniformly bounded.

Then, there exist constants \(K>0\), \(C>0\) such that for all \(t>0\),
\[
\mathbb P\!\left\{
\|\hat g - g^\star\|_{2}^2
\;\ge\;
\inf_{g\in\mathcal G}\|g - g^\star\|_{2}^2
\;+\; K\Bigl(\hat r_{\mathcal G}(n)^2 + \tfrac{t}{n}\Bigr)
\right\}
\;\le\; C e^{-t},
\]
where
\[
\hat r_{\mathcal G}(n)
\;:=\;\inf\Bigl\{r>0:\ \hat{\mathfrak{R}}_n(\mathcal G,r)\ \lesssim\ r^2\Bigr\},
\qquad
\hat{\mathfrak{R}}_n(\mathcal G,r)
:= \E_{\varepsilon}\!\left[
\sup_{\substack{g,h\in\mathcal G:\,\|g-h\|_2 \le r}}
\frac{1}{n}\sum_{i=1}^n \varepsilon_i\{g-h\}(X_i)
\right].
\]
\end{theorem}

We use the following corollary in the one-step regression analysis.

\begin{corollary}[PAC form; well-specified]
\label{cor::excessLS}
Under the conditions of Theorem~\ref{thm:ls-erm-rad} and assuming
\(g^\star \in \mathcal G\), for any \(\delta\in(0,1)\), with probability at
least \(1-\delta\),
\[
\|\hat g - g^\star\|_{2}^2
\;\le\;
K\Bigl(\hat r_{\mathcal G}(n)^2 + \tfrac{1}{n}\log\tfrac{1}{\delta}\Bigr).
\]
\end{corollary}

\begin{proof}
By Theorem~\ref{thm:ls-erm-rad}, for all \(t>0\),
\[
\Pr\!\left\{
\|\hat g-g^\star\|_{2}^2
\;\ge\; \inf_{g\in\mathcal{G}}\|g-g^\star\|_{2}^2
+ K\!\left(\hat r_{\mathcal{G}}(n)^2+\tfrac{t}{n}\right)
\right\}\le C e^{-t}.
\]
Under \(g^\star\in\mathcal{G}\), the infimum is \(0\). Set
\(t=\log(C/\delta)\) so that \(Ce^{-t}=\delta\). Then with probability at
least \(1-\delta\),
\[
\|\hat g-g^\star\|_{2}^2
\;\le\; K\!\left(\hat r_{\mathcal{G}}(n)^2+\tfrac{1}{n}\log\tfrac{C}{\delta}\right)
\;\le\; K'\!\left(\hat r_{\mathcal{G}}(n)^2+\tfrac{1}{n}\log\tfrac{1}{\delta}\right),
\]
absorbing \(\log C\) into \(K'\).
\end{proof}
